\begin{exercice}\label{exo1-5} (\mortelexo)

Nous avons vu en cours qu'un bon indicateur de localisation $c$ doit minimiser l' erreur
$e$ commise lorsque l'on r\'esume tout l' \'echantillon $(x_i)$ \`a l'aide de la seule valeur $c$.
Le calcul pr\'ecis de l'erreur $e$ d\'epend du choix de la distance qu'on souhaite utiliser.
Trouver la valeur de $c$ qui minimise $e$ pour les expressions suivantes de $e$:
\begin{enumerate}

\item {\it "Norme $L^1$"}:

\[ 
e = \frac1n \sum_{i=1}^n | x_i - c |.
\]

\item {\it "Norme $L^\infty$"}:

\[ 
e = \frac1n \max_{i=1,\ldots,n} | x_i - c |.
\]

\end{enumerate}

%\corrref{TD1_3}

\end{exercice}
